\chapter{拓扑学：橡皮泥眼中什么才算同一种形状}

\section{一张不像地图的地图}

你大概坐过地铁，也一定见过站台里那张花花绿绿的线路图：横平竖直的线条，均匀排布的圆点，整整齐齐的转角。但如果你把这张图和真实的城市地图叠在一起，会发现它们根本对不上——图纸上相距两厘米的两站，实际上可能隔了五公里；图纸上一条笔直的线路，实际上在地底下绕了好几个弯。这张地图在长度、角度、面积上全都撒了谎。

可奇怪的是，没有一个乘客因此迷路。一九三三年，绘图员哈利·贝克为伦敦地铁画出第一版这样的线路图时，地铁公司一开始还犹豫：把地图画得这么“不真实”，乘客能接受吗？结果事实证明，这张图好用到全世界所有地铁系统都照抄了它。为什么？因为坐地铁的人只关心三件事：我现在在哪一站，要去哪一站，中间在哪里换乘。至于地下隧道究竟拐了几度的弯、两站之间到底是八百米还是一千二百米，统统无所谓。贝克看透了这一点：他扔掉了一切会干扰视线的真实细节，只保留了**谁和谁相连**这一种信息。

这就引出了一个深刻的问题：一张地图可以扭曲得面目全非，却仍然“正确地”回答了某些问题。那么，一个图形身上到底有哪些性质，是随便怎么拉伸、弯曲、压缩都不会改变的？又有哪些性质，轻轻一捏就没了？数学家把研究前一类性质的学问叫作**拓扑学**。它有一个著名的绰号：橡皮泥几何学。在拓扑学家眼里，图形是用永远拉不断的橡皮泥做的——你可以随意揉捏它，只是不许撕开，也不许把两处粘在一起。被这样揉来揉去之后还纹丝不动的那些性质，就是拓扑学要研究的“真正的形状”。

本章我们就来当一回揉橡皮泥的人。我们会看到：一个三百年前的散步问题怎样催生了这门学科；咖啡杯和甜甜圈为什么在拓扑学家眼里是同一个东西；一张纸条拧半圈之后为什么会失去“正反面”；以及这些听起来像游戏的想法，怎样用在网络、地图和机器人身上。

\section{七座桥，一条走不出来的路线}

故事要从一座真实存在过的城市讲起。十八世纪的哥尼斯堡（今天俄罗斯的加里宁格勒）有一条河穿城而过，河中有两座岛，河两岸与两座岛之间一共架着七座桥。市民们茶余饭后有一个消遣问题，据说在酒馆里争论了很多年：

\begin{quote}
能不能从家里出发散步，把七座桥每座都恰好走一遍，最后回到家里？
\end{quote}

先别往下看，自己凭直觉猜一猜：能，还是不能？很多人第一反应是“能，多试几次总能找到路线”。也有人猜“不能”，但说不出理由。这正是这个问题迷人的地方：它听起来像一道脑筋急转弯，实际上谁也给不出让人信服的答案。城里的聪明人拿着地图一遍遍试，画出一条又一条路线，总是在最后一两座桥那里卡住——要么某座桥被迫走了两遍，要么走到一座岛上发现没有新桥可走了。

这种“试到猴年马月”的办法有一个致命的缺陷：就算你再试一万次都失败，也不能证明第一万零一次不会成功。要证明“办不到”，需要的不是更多的尝试，而是一种全新的眼光。一七三六年，当时住在圣彼得堡的数学家欧拉听说了这个问题。他先做了一件在当时看来很不“数学”的事：把地图扔掉了。

\section{把地图揉成一团之后，剩下什么}

欧拉的思路是：河岸有多大、岛是什么形状、桥有多长、路怎么拐——这些统统不重要。散步者真正经历的，只是“在某块陆地上”以及“从这块陆地经过某座桥到了另一块陆地”。于是他把四块陆地（河的南北两岸加两座岛）各自缩成一个点，把七座桥各自缩成连接两个点的一条线。得到的图丑极了，就四个点七条线，小学生三秒钟就能画完：

\begin{center}
\begin{tikzpicture}[scale=1.1]
  \node[circle,draw,fill=black,inner sep=2pt,label=above:{北岸 $A$}] (A) at (0,2) {};
  \node[circle,draw,fill=black,inner sep=2pt,label=below:{南岸 $B$}] (B) at (0,0) {};
  \node[circle,draw,fill=black,inner sep=2pt,label=left:{岛 $C$}] (C) at (-2.5,1) {};
  \node[circle,draw,fill=black,inner sep=2pt,label=right:{岛 $D$}] (D) at (2.5,1) {};
  \draw (A) to[bend left=25] (B);
  \draw (A) to[bend right=25] (B);
  \draw (C) -- (A);
  \draw (C) to[bend left=25] (A);
  \draw (C) -- (B);
  \draw (D) -- (A);
  \draw (D) -- (B);
\end{tikzpicture}
\end{center}

原来的问题在这张丑图上变成了：能不能从某个点出发，**每条线恰好描一遍**，最后回到出发点？这不就是我们小时候玩的“一笔画”吗！

请注意这一步的惊人之处。欧拉没有去解原来的问题，而是先证明了：原来的问题和这个一笔画问题是**同一个问题**。他把大地图揉成了一团小毛线，几何信息（长度、面积、形状）全都扔掉了，可问题所需要的全部信息居然一点没少。这种“只留下连接关系”的压缩手法，就是拓扑学的基本功。扔掉的东西越多，剩下的结构越显眼。

在继续之前，我们需要给这堆“点和线”起几个正式的名字，否则后面的话说不清楚。

\begin{dingyi}[图、顶点、边、度]
由若干个点（叫作\textbf{顶点}）和连接它们的若干条线（叫作\textbf{边}）组成的整体叫作一个\textbf{图}。一个顶点上连着几条边，这个数叫作该顶点的\textbf{度}。度为奇数的顶点叫\textbf{奇点}，度为偶数的叫\textbf{偶点}。如果从一个顶点出发沿着边能走到任何一个别的顶点，就称这个图是\textbf{连通}的。
\end{dingyi}

数数看，上面七桥图的四个顶点各是几度：$A$ 连着 $5$ 条边，$B$ 连着 $3$ 条边，$C$ 连着 $3$ 条边，$D$ 连着 $3$ 条边。四个顶点全是奇点。请记住这个事实，它马上就是破案的关键。

\begin{shuyu}{拓扑等价（同胚）}
直观解释：两个图形由橡皮泥做成，不用撕开、不用粘合，只靠连续地拉伸、弯曲、压缩就能从一个变成另一个，那么它们就是“同一种形状”。\\
正式定义：两个图形之间存在一个连续的一一对应，而且反过来的对应也连续，就称它们\textbf{同胚}。\\
例子：圆和正方形同胚（把圆周像皮筋一样撑成方框即可）；但圆和数字“8”不同胚——“8”中间那个交叉点拆不开。
\end{shuyu}

\section{破案：数一数为奇数的“路口”}

现在来解决一笔画问题。想象你拿着笔在某张图上走，经过某个顶点时，你必然是“从一条边进来，再从另一条边出去”——一进一出，成对消耗这个顶点的边。所以一个中途经过（既非起点也非终点）的顶点，它身上的边必须能配成一对一对的“进出口”，也就是说，它的度必须是偶数。

那么起点和终点呢？如果散步要回到出发点，起点被“出去一次、最后回来一次”，中间每经过一次也是一进一出，所以它的度同样是偶数。于是整个图里一个奇点都不许有。如果不要求回到出发点，那么起点“净出去一次”、终点“净进来一次”，允许它们是奇点——但也只允许这两个。这就是欧拉的发现：

\begin{dingli}[欧拉一笔画定理，1736]
一个连通的图：
\begin{enumerate}
  \item 存在一条经过每条边恰好一次并回到起点的路线，当且仅当图中\textbf{没有}奇点；
  \item 存在一条经过每条边恰好一次（不必回到起点）的路线，当且仅当图中\textbf{恰好有两个}奇点（此时它们必须充当起点和终点）。
\end{enumerate}
\end{dingli}

\begin{proof}
先证“必须有这些条件”。沿任何一条这样的路线走，每当中途经过一个顶点，都要消耗两条边（进一条、出一条），所以非起点终点的顶点必为偶点。若路线回到起点，起点消耗的两条边也成对，故起点也是偶点，图中没有奇点；若路线不回到起点，起点和终点各多出一条不配对用的边，它们恰是仅有的两个奇点。

再证“这些条件够用”。以没有奇点的情形为例。从任一顶点出发，沿着没走过的边随便走。由于每个顶点的边都成对，你走进任何一个顶点时，只要不是回到起点，就一定还有没走过的边可以出去——所以你唯一可能被“堵死”的地方就是起点。走到堵住为止，你必然回到了起点，得到一条回路。如果此时还有没走过的边，由于图是连通的，这条回路上一定有某个顶点连着没走过的边；从这个顶点出发再“随便走”一圈，同理又会回到它自己，把这新的一圈插进原来的回路里，得到一条更长的回路。如此反复，直到边被用完，就得到了所求的路线。恰好有两个奇点的情形类似：从一个奇点出发，最终必被堵在另一个奇点处，再用同样的方法插入剩余的回路即可。
\end{proof}

现在回头看七桥问题：四个顶点全是奇点，奇点个数是 $4$，既不是 $0$ 也不是 $2$。所以散步路线**不存在**——不是“暂时没找到”，而是“永远找不到”。欧拉没有多走一步路，只是数了数每个路口有几条桥，就宣判了全城人多年争论的死刑。据说哥尼斯堡后来建了第八座桥，奇点变成了两个，散步者终于可以走完所有桥（只是回不了家）。数学先预测了可能性，工程随后跟上，这是个颇有讽刺意味的后续。

\begin{toolbox}
本章到目前为止添置的新工具：
\begin{itemize}
  \item \textbf{压缩法}：把地图揉成“点加线”，只保留连接关系；
  \item \textbf{图}（顶点、边、度）和\textbf{连通性}：描述“谁和谁连着”的最小语言；
  \item \textbf{奇偶配对}：一进一出必成对，奇偶性是推理的发动机；
  \item \textbf{同胚}：判断“橡皮泥意义下同一种形状”的标准。
\end{itemize}
\end{toolbox}

\section{字母的分类游戏：什么是拓扑性质}

有了“同胚”这把尺子，我们可以玩一个分类游戏。把印刷体大写字母想象成橡皮泥做的细线：A、B、C、D、E……哪些是“同一种形状”？

C、I、J、L、M、N、S、U、V、W、Z 是一类：它们都是“一条线”，可以拉直成一根线段。A、D、O、P、Q、R 是一类：它们身上都有**一个洞**（A 的三角形孔、O 的圆孔，在橡皮泥眼里没区别）。B 独占一类：它有**两个洞**。E、F、T、Y 是一类：没有洞，但有**一个三叉路口**。K 和 X 是带四叉路口的一类。H 比较特殊，有两个三叉路口。

这个游戏告诉我们，拓扑学关心的是哪些性质呢？数一数上面的分类依据，无非就是这么几样：

\begin{itemize}
  \item \textbf{连通性}：图形是一整块，还是分成了好几块？（字母 i 上面的点算不算连着它的身子，就是个连通性问题。）
  \item \textbf{洞的数量}：有几个独立的“圈”绕不过去？
  \item \textbf{交叉与端点}：线在哪里分叉、在哪里终止？
  \item \textbf{边界}：图形有没有“边缘”，边缘有几条？
\end{itemize}

而长度、角度、面积、曲直，统统不在此列。拓扑学家有一句著名的自嘲：拓扑学家是分不清咖啡杯和甜甜圈的人。为什么？想象一个橡皮泥做的咖啡杯，杯身上有个把手。你把杯身凹陷进去、捏圆，只留下那个把手变成圆环——它就成了一只甜甜圈。全程没有撕开、没有粘合，杯子的洞（把手的孔）就是甜甜圈的洞，一个不多一个不少。反过来，甜甜圈永远变不成一个**没有洞的球**：球上的任何橡皮圈都能收缩成一点，而套在甜甜圈环上的圈怎么缩都会卡在环上。洞的数量，是把图形分成不同“物种”的第一道关卡。

\begin{fanli}
“两个图形看起来差得很远，所以它们拓扑上不同”——错。圆和正方形看起来差远了，但它们同胚；反过来，“两个图形都有同样的面积，所以拓扑上相同”也错：一个圆盘和一个挖了孔的圆环面积可以一样，但前者没洞、后者有一个洞，橡皮泥捏不出那个孔。判断同胚靠的是能否连续变形，不是看长相，也不是量大小。
\end{fanli}

\section{会数洞的数字：欧拉示性数}

洞是个好东西，但“数洞”有时说不清：一个打了三个孔的奶酪块，到底算几个洞？数学家需要一个更机械、更不会吵架的办法来记录图形“绕来绕去”的程度。办法出乎意料地朴素：把图形切碎成小块，数三种零件的个数，做一道加减法。

拿凸多面体开刀。一个立方体有 $8$ 个顶点、$12$ 条棱、$6$ 个面，算一算：
\[
V - E + F = 8 - 12 + 6 = 2.
\]
换一个，正四面体：$V=4$、$E=6$、$F=4$，算出来 $4-6+4=2$。再换一个，正十二面体：$V=20$、$E=30$、$F=12$，还是 $20-30+12=2$。有没有觉得不对劲？零件个数天差地别，加减完居然都是 $2$。欧拉（对，又是他，这个人几乎出现在数学的每一个角落）证明了这不是巧合：

\begin{dingli}[多面体的欧拉公式]
任何可以连续变形为球面的多面体（或画在球面上的连通图），都满足
\[
V - E + F = 2,
\]
其中 $V$、$E$、$F$ 分别是顶点数、棱数、面数。
\end{dingli}

\begin{proof}[证明大意]
把多面体想象成空心的橡皮气球，剪掉一个面，把剩下的部分摊平在桌面上，得到一个由顶点和边构成的连通图，它把平面分成若干区域，区域数比原来的面数少 $1$（被剪掉的那个面变成了图形外面的“无限区域”）。我们只要证摊平后的图满足 $V-E+F_{\text{区域}}=1$。现在做两种“拆除手术”：若图中有一个圈，删掉圈上的一条边，$E$ 少 $1$，圈两侧的两个区域合并，$F_{\text{区域}}$ 也少 $1$，$V-E+F_{\text{区域}}$ 不变；若没有圈了，图是一棵没有回路的“树”，此时删掉一个末端顶点连同连着它的那条边，$V$ 少 $1$、$E$ 少 $1$，$V-E+F_{\text{区域}}$ 还是不变。反复拆除，最后只剩一个孤零零的顶点：$V=1$、$E=0$、$F_{\text{区域}}=0$，数值为 $1$。既然手术从不改变这个数值，原来的图也一定等于 $1$，还原回多面体就是 $2$。
\end{proof}

这个数 $V-E+F$ 叫作\textbf{欧拉示性数}，通常记作希腊字母 $\chi$。它的神奇之处在于：你随便怎么切分这个曲面——切大块切小块、用直边用弯边——算出来的 $\chi$ 都一样。它是曲面的“指纹”。

而指纹能识别物种。把同样的游戏搬到甜甜圈（圆环面）上：在环面上画一个连通的图，把环面划分成若干块区域，你会算出 $V-E+F=0$。打两个洞的“双环面”（想象两节连在一起的游泳圈）算出来是 $-2$。规律是：每多一个洞，$\chi$ 就减 $2$。如果用 $g$ 表示洞的数量（数学里叫作\textbf{亏格}），那么
\[
\chi = 2 - 2g.
\]
球面 $g=0$，$\chi=2$；甜甜圈 $g=1$，$\chi=0$；双环面 $g=2$，$\chi=-2$。于是“数洞”这个说不清的任务，变成了一道谁算都一样的算术题：把曲面切成网格，数一数，减一减，就知道它有几个洞。咖啡杯能变成甜甜圈，正是因为它们的 $\chi$ 都是 $0$；它们都变不成球，因为 $\chi=0$ 永远不等于 $\chi=2$。

\begin{shuyu}{欧拉示性数}
直观解释：给曲面做一次“人口普查”，顶点数减边数加面数，是曲面在连续变形下不改写的身份证号。\\
正式定义：把曲面划分成 $V$ 个顶点、$E$ 条边、$F$ 个面，$\chi = V - E + F$；它与划分方式无关。\\
例子：立方体 $8-12+6=2$，所以立方体和球面同类；环面上的网格算出 $0$，比球面少 $2$，正好对应多出来的一个洞。
\end{shuyu}

\section{拧半圈的世界：莫比乌斯带与克莱因瓶}

到目前为止，我们的橡皮泥游戏还保留着一条没说出口的规矩：曲面有里外两面。气球有内壁和外壁，纸片有正面和反面。可一八五八年，数学家莫比乌斯发现了一种只有**一面**的曲面，制作方法简单到荒谬：取一张长纸条，把一端旋转半圈，再和另一端粘起来。

做一条试试（本章后面的实验环节会带你做）。在纸带上用彩笔沿着中线画线，笔尖不离开纸面，画着画着你会发现：笔回到了出发点，而你已经把“两面”全涂满了——这条纸带根本没有“另一面”！沿中线剪开，它不像直觉预期的那样分成两个环，而是变成一个拧了一整圈的、更长的大环。一只在莫比乌斯带上爬行的蚂蚁，可以不翻越任何边缘就爬遍整条带子。

\begin{shiyan}
找一张废纸，裁下一条约 $3$ 厘米宽、$25$ 厘米长的纸条。
\begin{enumerate}
  \item 先把两端直接粘起来，得到一个普通的圆环。用彩笔沿环的中线画一圈，笔只涂到了“一面”。沿中线剪开，圆环分成两个环。
  \item 重新做一条，这次粘之前把一端拧半圈。沿中线画线——这次笔迹“神奇地”布满了整条带子。沿中线剪开，数一数得到的环拧了几圈。
  \item 再拧一整圈（$360^{\circ}$）粘起来，重复画线和剪开，结果和前两次各有什么不同？
  \item 记录：拧 $0$ 圈、半圈、一圈、一圈半时，剪开后分别得到几个环、各拧了几圈。你能猜出规律吗？
\end{enumerate}
\end{shiyan}

莫比乌斯带还有一条**边界**（就是你没剪之前那条长长的纸边，整条带子的边缘居然是连通的一条线）。数学家接着问：能不能造出一种只有一面、而且连边界都没有的封闭曲面？可以，这就是\textbf{克莱因瓶}。想象把一个瓶子的瓶底打个洞，把瓶颈拉长、弯过来、穿过瓶壁，和瓶底的洞粘在一起。做成之后，瓶子的“里面”和“外面”是同一片曲面——沿着瓶壁走，你可以从外壁不知不觉走到内壁，无须穿过任何边缘。一只蚂蚁能爬遍整个克莱因瓶。

但这里有个不得不坦白的事实：克莱因瓶在我们生活的三维空间里**造不出来**。刚才的描述里有一句“穿过瓶壁”——在三维世界中，瓶颈想接到瓶底，就必须穿透瓶身，这就违反了橡皮泥“不许穿孔”的家规。只有在四维空间里，瓶颈可以从“第四个方向”绕过去，不必碰瓶壁。你在工艺品店里买到的“克莱因瓶”玻璃摆件，全都有一处自交，那是对这个四维居民的三维投影，就像甜甜圈的影子落在墙上可能自带交叉一样。莫比乌斯带告诉我们“面数”不是天经地义的，克莱因瓶则告诉我们：有些形状大到我们的世界装不下，却照样可以被数学精确地描述。

\begin{lianjie}
还记得本书前面把正方形、旋转和翻折编成“对称的运算表”吗？那里我们把“看起来一样”升级成了“结构一样”。本章做的是同一种升级，只是换了尺子：第 6 章问“做什么动作后图形不变”，本章问“怎么揉捏后性质不变”。而本章的“图”与第 14 章的图论是同一套语言——七桥问题正是图论和拓扑学共同的出生证明。下一章你将看到，当橡皮泥换回直尺和量角器，故事又会怎样翻出新的一页。
\end{lianjie}

\section{换一个视角：网络、地图和机器人}

拓扑学听起来像在玩橡皮泥和剪纸，但它的“保留连接、扔掉细节”的哲学，恰恰是现代世界运转的底层逻辑。

\textbf{网络。}互联网有几亿台设备，城市交通网有成千上万个路口，电网有数不清的节点。工程师关心的第一件事往往不是“这根电缆多长”，而是拓扑问题：这个网络连通吗？断掉某个节点，网络会不会裂成两半？裂成两半时各有多大？这与本章开头判断“图形是几整块”是同一种思考。社交网络的“六度分隔”现象——任意两个人之间平均只隔着大约六层朋友关系——说的正是巨大网络在连接关系上出人意料地“紧”。

\textbf{地图。}给世界地图上色，要求相邻国家颜色不同，最少要几种颜色？注意这个问题只跟“谁和谁接壤”有关，跟国家的面积、形状、海岸线弯曲程度毫无关系——这是纯拓扑问题。人们猜了上百年“四种颜色足够”，这就是著名的四色猜想。一九七六年，数学家把它化归为一千多种基本图形的逐一检查，借助计算机完成了证明，成为第一个“计算机参与证明”的大定理，至今仍有数学家争论：一个人无法亲手核验的证明，算不算证明？

\textbf{机器人。}一个机械臂有几个关节，每个关节可以转动一个角度。机械臂的每一个姿势，都对应着由这些角度组成的一个“形状空间”里的一个点；规划机械臂的运动，就是在这个空间里找一条不撞墙的连续路径。奇妙的是，这个空间的拓扑形状本身就有讲究：一个能整圈旋转的关节贡献一个“圆圈”，两个独立旋转的关节贡献一个“环面”——正是本章那只甜甜圈！机器人工程师规划路径时，实际上是在甜甜圈、球面或更怪的曲面上走路。路径规划里的很多困难，比如“怎么转都绕不过去的障碍”，本质上就是空间里的“洞”在作怪。

\begin{zhu}
以上应用里，拓扑负责回答“能不能”：能不能连通、能不能绕过、能不能四色。至于“多快、多长、多省钱”，那要把长度和角度请回来，交给别的数学分支。分清一个问题属于“能还是否”还是“好还是坏”，是解题的第一步。
\end{zhu}

\begin{pandeng}
想继续往前走，可以记住这几个关键词：\textbf{基本群}（把曲面上“缩不回去的圈”编成一种可以运算的代数结构，洞的数量由此变成严格的定理）；\textbf{同调论}（欧拉示性数的远房亲戚，能数出任意维度的“洞”）；\textbf{纽结理论}（研究三维空间里的绳结何时能解开，你的鞋带和 DNA 都在里面）；\textbf{庞加莱猜想}（三维世界版的“没有洞的封闭形状必是球面”，直到二〇〇三年才被证明，是七大千禧年难题中第一个被攻克的）。
\end{pandeng}

\section*{思考题}

\textbf{观察题}

\begin{sikao}
  \item 把阿拉伯数字 $0$ 到 $9$ 看作橡皮泥细线（手写体，不带交叉），按拓扑形状分类。哪几个数字属于同一类？\textbf{提示}：先数每个数字有几个洞、几个端点。
  \item 一件长袖衬衫（没有系扣子）平铺后，按拓扑形状算，它和一个打了几个洞的圆盘同类？\textbf{提示}：数一数身体伸进去之后，一共有几个“出口”，想想每个出口对应几个洞。
\end{sikao}

\textbf{证明题}

\begin{sikao}
  \item 证明：任何一个图里，奇点的个数一定是偶数。\textbf{提示}：把所有顶点的度加起来，每一条边被数了几次？这个总和是奇数还是偶数？
  \item 用 $V-E+F$ 检验：在一个球面上画一个“经纬网”——比如 $4$ 条经线、$2$ 条纬线（赤道和一条），南北极点也算顶点。数一数顶点、边、区域，验证欧拉公式。\textbf{提示}：先想清楚两条经线在极点相交算不算顶点，每条纬线被经线切成几段。
\end{sikao}

\textbf{挑战题}

\begin{sikao}
  \item 甜甜圈上最多能画几个两两都直接相邻的区域（每两个区域之间都有一条公共边界）？平面上的答案至多是 $4$（这正是四色定理里“四”的来源之一），甜甜圈上的答案比它大。\textbf{提示}：答案是 $7$。先试着画 $5$ 个两两相邻的区域，再利用“环可以绕过去”多塞两个。
  \item 把两张莫比乌斯带沿各自的边界粘起来，会得到本章介绍过的哪种曲面？\textbf{提示}：先在脑子里把一张莫比乌斯带“展开”成一个圆环带子；想想克莱因瓶能不能沿着某条线剪成两半。
\end{sikao}

\section*{通往下一章的门}

这一章我们过得相当潇洒：长度可以扔，角度可以扔，面积可以扔，只要连接关系和洞的数量还在，图形就算“没变”。橡皮泥的世界里，球面和土豆是同一种东西，茶杯和甜甜圈不分彼此。

可是地球人不会同意“地球和土豆一样”——导航软件要算距离，飞行员要知道从北京飞纽约为什么先往北飞，而不是沿纬线直着往东。这些问题一上来就要动用本章扔掉的东西：长度、角度、弯曲程度。球面上的三角形内角和居然大于 $180^{\circ}$，而甜甜圈上的直线会自己缠上自己——这些现象拓扑学一律看不见，因为橡皮泥一捏就全没了。

如果我们把橡皮泥换回一块硬一点的、有弧度的曲面，不许拉伸但允许在上面滑行，那么“最直的路”“最短的距离”“弯得有多厉害”这些词，就需要一套全新的语言来定义。曲面在每个点附近怎样近似成一张平面？一只蒙着眼睛的蚂蚁，能不能只凭脚下的感觉判断自己住在球面上还是平面上？下一章，我们走进微分几何：弯曲空间中的距离和方向。
